% Figure dq: generated by figures/rse_fig_dq.py; placed where the text first discusses it.
% Construction of the degraded-elevation-model control: Supplementary section rse_si_display.tex.
\begin{figure}[!tbp]
  \centering
  \includegraphics[width=\textwidth]{../figures/out/rse_fig_dq.pdf}
  \caption{\textbf{Grid error against observation error, by sector and grid.}
  For each grid and sector, the ratio $D/Q$ at every level at which both can be measured ($n = \metric{h1.n_evaluable.shape.gebco2026.west}$--\metric{h1.n_evaluable.shape.gebco2026.interfluve} levels per grid and sector).
  $D$ is the grid's median shoreline displacement at the level, after its area-fitted vertical offset; $Q$ is the observation noise floor at the level, the \metric{params.err_pct_hi}th percentile of the distances between the shorelines of two interleaved halves of the acquisitions.
  (a)~The four checked sectors; (b)~the three sectors whose observations are biased against ICESat-2 or have no test block.
  Dots: one level each. Boxes: \metric{params.iqr_lo_pct}th to \metric{params.iqr_hi_pct}th percentile over levels. Large markers: the median; filled where the grid fails the shoreline test in the sector ($D > Q$ at more than half of the levels), open where it passes.
  Grey, to the right of each grid: its degraded-elevation-model control (Section~\ref{sec:sep}), $D_\mathrm{control}/Q$, at the levels where both can be measured; the two GEBCO grids share one control.
  Shaded: $D$ below $Q$. Ratios of zero are drawn on the lower edge. Logarithmic axis.}
  \label{fig:rse_dq}
\end{figure}
